\chapter{Metode și Configurare Experimentală}
\label{ch:methodology}

% This chapter explains every model and every shared experimental choice
% (preprocessing, hyperparameters, evaluation protocol) that applies across
% all three datasets. Dataset-specific preprocessing details that don't
% generalize (e.g. exact feature counts, window lengths) can either stay
% here at a high level or move into the per-dataset Results subsections --
% decide per case once the prose is drafted.

Hiperparametrii și alegerile legate de antrenare, per model (optimizator, epoci, lățimi de straturi și restul detaliilor de arhitectură, câte un tabel per model pentru cele patru variante de set de date/caracteristici) sunt prezentate în Anexa~\ref{sec:architecture}.

\section{Preprocesare}
\label{sec:preprocessing}

Două principii se aplică la toate cele trei pipeline-uri: (1) fiecare model este antrenat exclusiv pe eșantioane benigne/normale (\emph{novelty detection}) -- partiția de antrenare nu conține niciodată o anomalie, iar scalerele/normalizatoarele sunt \emph{potrivite} doar pe antrenare, apoi aplicate neschimbate pe testare, evitând data leakage; 
(2) partiția de testare conține întotdeauna ambele clase, evaluate cu protocolul din \S\ref{sec:eval-protocol}. Multe dintre detaliile considerate neesențiale pentru înțelegerea textului au fost mutate în Anexa~\ref{sec:preprocessing-details}.

\subsection{Preprocesare Tabelară (Credit Card)}

\texttt{Time} este eliminat (nu reprezintă informație semnificativă). Rândurile benigne (\texttt{Class=0}) sunt împărțite 80/20 în antrenare/test-benign, iar setul de testare adaugă \emph{toate} cele 492 de rânduri de fraudă. Scalarea este asimetrică: \texttt{V1--V28} rămân nescalate (deja caracteristici PCA), iar \texttt{Amount} este scalat cu un \texttt{RobustScaler}, ales pentru robustețea sa la coada dreaptă accentuată. Pași exacți: Anexa~\ref{sec:preprocessing-details}.

\subsection{Crearea de Ferestre pentru Serii de Timp (Yahoo S5)}

Fiecare serie este împărțită temporal 70/30 (\texttt{TRAIN\_FRAC=0.70}, fără amestecare aleatoare, pentru a nu lăsa informație din viitor să se scurgă în antrenare), scalată \emph{per serie} cu un \texttt{MinMaxScaler} potrivit doar pe punctele neanormale de antrenare, apoi transformată în ferestre de lungime \texttt{WINDOW=64} (pas 4 la antrenare, pas 1 la testare). Eticheta unei ferestre urmează regula \enquote{oricare}: pozitivă dacă cel puțin un pas este anormal. Detalii complete, plus comparația metodologică cu referința DeepAnt~\parencite{munir2019deepant}: Anexa~\ref{sec:preprocessing-details}.

\subsection{Preprocesarea Fluxurilor de Rețea (CIC-IDS2017)}

Există trei implementări, toate având la bază aceiași pași de curățare (eliminarea coloanelor de identificare, codificare one-hot pentru \texttt{Protocol}, imputare cu mediana, \texttt{MinMaxScaler} potrivit pe antrenare, împărțire pe zile: luni antrenare, marți--vineri testare), dar diferind prin unitatea de observație: \textbf{Plat} (un rând = un flux, 77 caracteristici), \textbf{Temporal} (un rând = o fereastră de \texttt{WINDOW=16} fluxuri consecutive de la aceeași adresă IP sursă) și \textbf{Context de gazdă} (fluxuri individuale suplimentate cu 9 caracteristici cauzale dintr-o fereastră glisantă de 60s per IP gazdă -- număr de fluxuri, porturi unice, rata de diversitate a porturilor etc. -- urmând adaptarea nesupervizată a lui \textcite{raskovalov2024nids}, 86 caracteristici în total). Gruparea se face după IP sursă, întrucât comportamentul la nivel de gazdă caracterizează atacuri precum PortScan sau DoS mai bine decât o singură conexiune. Detalii complete: Anexa~\ref{sec:preprocessing-details}.

\section{Protocolul de Evaluare}
\label{sec:eval-protocol}

La testare, fiecare model produce un scor scalar de anomalie per eșantion -- niciodată o etichetă directă -- iar această etichetă este obținută prin aplicarea unui prag asupra scorului, urmând procedura descrisă mai jos.

\subsection{Metrici Raportate}

Lista completă a metricilor raportate de fiecare rulare, precum și rezultatele complete pe toate aceste metrici pentru fiecare model și set de date, sunt prezentate în Anexa~\ref{sec:eval-metrics}.

PR-AUC (Precizia Medie) este tratată drept metrica principală pentru compararea modelelor, în locul F1 sau ROC-AUC, urmând raționamentul lui \textcite{kim2022rigorous}: prestabilirea unui prag de decizie fără acces la setul de testare nu este practică, iar metodele existente raportează în schimb F1 la orice prag care se întâmplă să îl maximizeze, ceea ce face rezultatele dependente de prag prin construcție. Aceștia concluzionează că \enquote{metrici suplimentare cu dependență redusă, precum AUROC sau aria de sub curba precizie-recall (AUPR), vor ajuta la o evaluare riguroasă} -- justificarea adoptată aici pentru clasificarea modelelor după PR-AUC.

ROC-AUC este de asemenea raportat, mai mult pentru completitudine/comparabilitate cu literatura anterioară.

Pentru CIC-IDS2017, este raportată suplimentar Precizia Medie defalcată pe tip de atac, pentru a susține discuția per-tip-de-atac din \S\ref{ch:results}.

\subsection{Selecția Pragului}

Se alege un singur prag de decizie per model/set de date pentru metricile secundare, dependente de prag (F1, acuratețe echilibrată, matrice de confuzie): se baleiază toate pragurile de pe curba precizie-recall și se alege cel care maximizează $F_\beta$ cu $\beta=2$, ponderând recall-ul dublu față de precizie. Valoarea lui $\beta$ nu este reglată -- orice $\beta>1$ ar exprima aceeași preferință operațională rezonabilă, unde o anomalie ratată e de obicei mai costisitoare decât o alarmă falsă; $\beta=2$ e o alegere reprezentativă.

Acesta este un prag de tip \emph{oracol}, selectat folosind chiar curba precizie-recall a setului de testare, ceea ce nu ar fi mereu disponibil în realitate. Este acceptabil aici întrucât obiectul de studiu este comparația relativă între modele la cel mai bun punct de operare, iar PR-AUC nu depinde de selecția pragului -- cifrele F2/F1/acuratețe echilibrată la pragul oracol rămân secundare, ilustrative.

\subsection{Fără Protocolul \emph{Point Adjustment}}

Multe lucrări despre detecția anomaliilor în serii de timp aplică \emph{Ajustarea de Punct} (\emph{Point Adjustment}, PA): dacă orice punct dintr-un segment contiguu de anomalie este semnalat, întregul segment e considerat detectat.
\textcite{kim2022rigorous} arată că aceasta poate umfla F1 până aproape de 1 chiar și pentru scoruri aleatorii, atunci când segmentele de anomalie sunt lungi, iar F1 ajustat prin PA este aproape necorelat cu F1 neajustat (Pearson $r=-0.59$ pe SWaT, în studiul lor). Nicio ajustare PA nu este aplicată nicăieri în acest proiect: etichetele sunt fixate la preprocesare, iar predicțiile sunt evaluate exact așa cum rezultă din pragul aplicat, fără ajustare retroactivă.